\documentclass[12pt,a4paper]{article}

\usepackage{fontawesome5}
\providecommand{\faPencilAlt}{\faIcon{pencil-alt}}

\usepackage{styles/couleurs}
\usepackage{styles/mise_en_page}
\usepackage{styles/boites}
\usepackage{styles/maths}
\usepackage{styles/figures}

\usepackage{enumitem}
\usepackage{array}
\usepackage{tabularx}
\usepackage{amsmath}
\usepackage{amssymb}
\usepackage{tikz}
\usepackage{xparse}

\usetikzlibrary{calc,angles,quotes,arrows.meta}

\setlist[itemize]{leftmargin=1.4em,itemsep=0.15em,topsep=0.2em}
\setlist[enumerate]{leftmargin=*,itemsep=0.22em,topsep=0.25em}

\newcommand{\lignepointillee}{\noindent\dotfill\par\vspace{0.45em}}

\newcommand{\lignesreponse}[1]{%
\begin{tcolorbox}[
enhanced,colback=white,colframe=NeutreC!40,boxrule=0.5pt,arc=1.5mm,
left=3mm,right=3mm,top=2mm,bottom=1mm,before skip=2pt,after skip=2pt]
\foreach \i in {1,...,#1}{\lignepointillee}
\end{tcolorbox}}

\newcommand{\titreseq}[1]{%
\begin{center}
{\sffamily\LARGE\bfseries\color{TitreC}Séquence 6 -- Cosinus d'un angle aigu}\\[0.2em]
{\sffamily\large\bfseries\color{SavoirC}#1}
\end{center}}

\begin{document}

\titreseq{Activité 7 -- Mesurer une hauteur inaccessible}

Une observatrice se tient à \(30\) m d'un arbre.
Ses yeux sont à \(1{,}75\) m du sol.
La ligne horizontale passant par ses yeux rencontre le tronc en \(H\).
Elle voit le sommet \(A\) sous un angle de \(35^\circ\).

On suppose le sol horizontal et le tronc vertical.

\begin{center}
\begin{tikzpicture}[scale=0.13]
  \coordinate (S) at (0,12);
  \coordinate (H) at (30,12);
  \coordinate (A) at (30,33);
  \coordinate (P) at (0,0);
  \coordinate (T) at (30,0);

  \draw[line width=1pt] (-3,0)--(34,0);
  \draw[line width=1pt] (T)--(A);
  \draw[line width=0.9pt] (S)--(H);
  \draw[line width=0.9pt] (S)--(A);
  \draw[dashed] (P)--(S);
  \draw[dashed] (H)--(T);

  \draw (30,12)--(28.2,12)--(28.2,13.8);

  \node[left] at (S) {$S$};
  \node[right] at (H) {$H$};
  \node[above] at (A) {$A$};
  \node[below] at (P) {$P$};
  \node[below] at (T) {$T$};

  \node[below] at (15,0) {$PT=30$ m};
  \node[left] at (0,6) {$PS=1{,}75$ m};

  \draw[ActionC,line width=1.2pt] (4,12) arc[start angle=0,end angle=35,radius=4];
  \node[ActionC] at (5.3,14) {$35^\circ$};
\end{tikzpicture}
\end{center}

\section*{Partie A -- Comprendre la figure}

\begin{enumerate}
  \item Pourquoi le triangle \(SHA\) est-il rectangle en \(H\) ?
  \item Quelle longueur vaut \(30\) m dans ce triangle ?
  \item La hauteur cherchée est-elle directement \(AH\) ? Explique.
  \item Quelle est l'hypoténuse du triangle \(SHA\) ?
\end{enumerate}

\lignesreponse{5}

\section*{Partie B -- Première étape : calculer la ligne de visée}

Dans le triangle \(SHA\) rectangle en \(H\), écris le cosinus de l'angle
\(\widehat{HSA}=35^\circ\).

Déduis la longueur \(SA\), arrondie seulement pour l'affichage,
mais conserve la valeur complète dans la calculatrice pour la suite.

\lignesreponse{6}

\section*{Partie C -- Deuxième étape : calculer la hauteur au-dessus des yeux}

Utilise le théorème de Pythagore dans le triangle \(SHA\)
pour calculer \(AH\).

\lignesreponse{6}

\section*{Partie D -- Conclure}

Calcule la hauteur totale \(AT\) de l'arbre.

Donne un résultat arrondi au centimètre et rédige une phrase de conclusion.

\lignesreponse{5}

\section*{Partie E -- Contrôler le résultat}

\begin{enumerate}
  \item La ligne de visée \(SA\) doit-elle être plus grande ou plus petite que \(SH\) ?
  \item Pourquoi \(AH\) est-il inférieur à \(SA\) ?
  \item Que se passerait-il si l'on oubliait la hauteur des yeux ?
\end{enumerate}

\lignesreponse{4}

\aretenir{Dans un problème complexe, le cosinus ne donne pas toujours directement
la longueur finale. Il peut fournir une longueur intermédiaire,
puis un autre outil géométrique permet de poursuivre.}

\end{document}
